%------------------------------------------------------------------------------%

\usepackage{calculo_varias_variaveis-1}

\title {Diferenciais}

%------------------------------------------------------------------------------%
\begin{document}

\addtitle

%------------------------------------------------------------------------------%
\section{Diferenciais}

%------------------------------------------------------------------------------%
\begin{frame}{Diferenial em Uma Dimensão}
  Quando $x$ varia de $x_0$ para
  \tabto{55mm}\(x_0+\Delta x\)

  \pause
  \vspace{\baselineskip}
  A variação da função $f$ é
  \tabto{55mm}\(\Delta f = f(x_0 + \Delta x) - f(x_0)\)

  \pause
  \vspace{\baselineskip}
  A diferencial de $f$ é
  \tabto{55mm}\(df = f'(x_0) \Delta x\)
\end{frame}

%------------------------------------------------------------------------------%
\begin{frame}{Diferenial em Uma Dimensão}
  Quando $x$ varia de $x_0$ para
  \tabto{55mm}\(x_0+dx\)

  \vspace{\baselineskip}
  A variação da função $f$ é
  \tabto{55mm}\(\Delta f = f(x_0 + dx) - f(x_0)\)

  \vspace{\baselineskip}
  A diferencial de $f$ é
  \tabto{55mm}\(df = f'(x_0) dx\)
\end{frame}

%------------------------------------------------------------------------------%
\begin{frame}{Diferenial em Duas Dimensão}
  Quando $(x, y)$ varia de $(x_0, y_0)$ para
  \tabto{65mm}\((x_0+dx, y_0+dy)\)

  \pause
  \vspace{\baselineskip}
  A variação da função $f$ é
  \tabto{65mm}\(\Delta f = f(x_0+dx, y_0+dy) - f(x_0, y_0)\)

  \pause
  \vspace{\baselineskip}
  A diferencial de $f$ é
  \tabto{65mm}\(df = f_x(x_0, y_0) dx +  f_y(x_0, y_0) dy\)

  \pause
  \vspace{\baselineskip}
  As diferenciais $dx$ e $dy$ são variáveis independentes
  \[
    dx = x - x_0
    \qquad
    dy=y-y_0
  \]
\end{frame}

%------------------------------------------------------------------------------%
\begin{frame}{Definição}
  Se movemos de \((x_0, y_0)\) para um ponto próximo \((x_0+dx, y_0+dy)\)

  \pause
  \vspace{\baselineskip}
  A \emph{Diferencial Total de $f$} é
  \[
    df \;=\; f_x(x_0, y_0)\, dx \;+\;  f_y(x_0, y_0)\, dy
  \]
\end{frame}

%------------------------------------------------------------------------------%
\addtocounter{exemplo}{1}
\begin{frame}{Exemplo \theexemplo}
  Uma lata de formato cilíndrico foi projetada para ter um raio de 1 pol.
  e uma altura de 5 pol., mas o raio tem um erro de \(dr=0,03\) e a altura
  tem um erro de \(dh=-0,1\).

  \vspace{\baselineskip}
  Qual a variação resultante no volume da lata?
\end{frame}

\begin{frame}{Exemplo \theexemplo}
  \onslide<+->{Variação exata}
  \begin{align*}
    \onslide<+->{\Delta V}
    \onslide<+->{& = V(r_0+dr, h_0+dh)- V(r_0, h_0) \\}
    \onslide<+->{& = V(1+0,03, 5-0,1) - V(1, 5)\\}
    \onslide<+->{& = V(1,03, 4,9) - V(1, 5) \\}
    \onslide<+->{& = \pi(1,03)^2 \times 4,9 - \pi 1^2 \times 5 \\}
    \onslide<+->{& \approx 16,3313 - 15,7080 \\}
    \onslide<+->{& = 0,6233}
  \end{align*}
\end{frame}

\begin{frame}{Exemplo \theexemplo}
  \onslide<+->{Variação usando o diferencial}
  \begin{align*}
    \onslide<+->{\Delta V}
    \onslide<+->{& \approx dV = V_r(r_0, h_0) dr + V_h(r_0, h_0) dh \\}
    \onslide<+->{& = \left(2\pi r h\right)\evalat{(r_0, h_0)}{} dr
      + \left(\pi r^2\right)\evalat{(r_0, h_0)}{} dh \\}
    \onslide<+->{& = \left(2\pi r h\right)\evalat{(1, 5)}{} (0,03)
      + \left(\pi r^2\right)\evalat{(1, 5)}{} (-0,1) \\}
    \onslide<+->{& = \left(2\pi 1 \times 5\right)(0,03)
      - \left(\pi 1^2\right)(0,1) \\}
    \onslide<+->{& = 0,3\pi - 0,1\pi  \\}
    \onslide<+->{& = 0,2\pi \\}
    \onslide<+->{& \approx 0,6283}
  \end{align*}
\end{frame}

%------------------------------------------------------------------------------%
\section{Lista Mínima}

%------------------------------------------------------------------------------%
\begin{frame}{Lista Mínima}
  Cálculo Vol. 2 do Thomas 12\textsuperscript{a} ed. -- Seção 14.6
  \begin{enumerate}
    \item Estudar o texto da seção
    \item Resolver os exercícios: 
  \end{enumerate}

  \vfill
  Atenção: A prova é baseada no livro, não nas apresentações
\end{frame}

%------------------------------------------------------------------------------%
\end{document}
%------------------------------------------------------------------------------%
